\documentclass{article}
\usepackage{tikz}
\usepackage{tikz,lipsum,lmodern}
\usepackage[most]{tcolorbox}
\usepackage{amssymb}
\usepackage{amsmath}
\usepackage{soul}
\usepackage{lastpage}
\usepackage{fancyhdr}
\usepackage{sectsty}
\usepackage{siunitx}
\usetikzlibrary{calc}
\usepackage{xcolor}


\pagestyle{fancy}
\fancyhf{}
\fancyfoot[C]{{\huge\textbf{\textcolor{aquablue!50!white}{\thepage /}}}{\Large \textbf{\textcolor{cyan!70!white}{\pageref{LastPage}}}}}
\renewcommand{\headrulewidth}{0pt}

\definecolor{purewhite}{RGB}{255, 255, 255}
\definecolor{offwhite}{RGB}{220, 225, 230}
\definecolor{electricblue}{RGB}{0, 191, 255}
\definecolor{deepblue}{RGB}{0, 102, 204}
\definecolor{navyblue}{RGB}{0, 51, 102}
\definecolor{darkcharcoal}{RGB}{15, 15, 20}
\definecolor{ivory}{RGB}{250, 249, 246}
\definecolor{royalblue}{RGB}{65, 105, 225}     
\definecolor{aquablue}{RGB}{0, 255, 255}       
\definecolor{oceanblue}{RGB}{0, 119, 182}     
\definecolor{aquamarine}{RGB}{127, 255, 212}   
\definecolor{babyblue}{RGB}{137, 207, 240}     
\definecolor{turquoise}{RGB}{64, 224, 208}     
\definecolor{darkturquoise}{RGB}{0, 206, 209}

\color{ivory}
\pagecolor{black}


\begin{document}
\begin{titlepage}
   \begin{tikzpicture}[remember picture,overlay,inner sep=0,outer sep=0]
     \draw[blue,line width=4pt] ([xshift=-1.5cm,yshift=-2cm]current page.north east) coordinate (A)--([xshift=1.5cm,yshift=-2cm]current page.north west) coordinate(B)--([xshift=1.5cm,yshift=2cm]current page.south west) coordinate (C)--([xshift=-1.5cm,yshift=2cm]current page.south east) coordinate(D)--cycle;

     \draw[aquablue, line width=2pt]([yshift=0.5cm,xshift=-0.5cm]A)-- ([yshift=0.5cm,xshift=0.5cm]B)--
     ([yshift=-0.5cm,xshift=0.5cm]B) --([yshift=-0.5cm,xshift=-0.5cm]B)--([yshift=0.5cm,xshift=-0.5cm]C)--([yshift=0.5cm,xshift=0.5cm]C)--([yshift=-0.5cm,xshift=0.5cm]C)-- ([yshift=-0.5cm,xshift=-0.5cm]D)--([yshift=0.5cm,xshift=-0.5cm]D)--([yshift=0.5cm,xshift=0.5cm]D)--([yshift=-0.5cm,xshift=0.5cm]A)--([yshift=-0.5cm,xshift=-0.5cm]A)--([yshift=0.5cm,xshift=-0.5cm]A);


     \draw [aquablue,line width=2pt]([yshift=-0.3cm,xshift=0.3cm]A)-- ([yshift=-0.3cm,xshift=-0.3cm]B)--
     ([yshift=0.3cm,xshift=-0.3cm]B) --([yshift=0.3cm,xshift=0.3cm]B)--([yshift=-0.3cm,xshift=0.3cm]C)--([yshift=-0.3cm,xshift=-0.3cm]C)--([yshift=0.3cm,xshift=-0.3cm]C)-- ([yshift=0.3cm,xshift=0.3cm]D)--([yshift=-0.3cm,xshift=0.3cm]D)--([yshift=-0.3cm,xshift=-0.3cm]D)--([yshift=0.3cm,xshift=-0.3cm]A)--([yshift=0.3cm,xshift=0.3cm]A)--([yshift=-0.3cm,xshift=0.3cm]A);

   \end{tikzpicture}
\begin{center}
  \centerline{{\Huge \bfseries \textcolor{aquablue!60!blue}{International Physics Competition}}}
  \vspace{1cm}
  {\Large \bfseries \color{electricblue} Qualification Round 2026}\\[1cm]
  
  {\large \bfseries \color{cyan!30!white} Solution Document}\\[2cm]
  
  {\color{aquablue} \Huge $\diamond$}\\[2cm]
  
  {\huge \bfseries \color{aquablue} Md Mokhdum Amin Fahim}\\[1cm]
  {\Large \bfseries\color{aquamarine} Notre Dame College, Dhaka}\\[1cm]
  {\Large \color{aquamarine!50!aquablue}\bfseries Bangladesh}\\[2cm]
  {\color{aquablue} \Huge $\diamond$}\\[2cm]
  {\large \bfseries \color{cyan!70!white} Date of Submission : 1 November, 2026}\\[1cm]
  
  {\large \bfseries \color{babyblue} Number of Pages : \pageref{LastPage} }\\[1cm]
 
  
  \vspace{\fill}
\end{center}
\end{titlepage}

\newpage
\setcounter{page}{1}
\begin{tikzpicture}[remember picture,overlay,inner sep=0,outer sep=0]
     \draw[blue,line width=4pt] ([xshift=-1.5cm,yshift=-2cm]current page.north east) coordinate (A)--([xshift=1.5cm,yshift=-2cm]current page.north west) coordinate(B)--([xshift=1.5cm,yshift=2cm]current page.south west) coordinate (C)--([xshift=-1.5cm,yshift=2cm]current page.south east) coordinate(D)--cycle;

     \draw [aquablue,line width=2pt]([yshift=0.5cm,xshift=-0.5cm]A)-- ([yshift=0.5cm,xshift=0.5cm]B)--
     ([yshift=-0.5cm,xshift=0.5cm]B) --([yshift=-0.5cm,xshift=-0.5cm]B)--([yshift=0.5cm,xshift=-0.5cm]C)--([yshift=0.5cm,xshift=0.5cm]C)--([yshift=-0.5cm,xshift=0.5cm]C)-- ([yshift=-0.5cm,xshift=-0.5cm]D)--([yshift=0.5cm,xshift=-0.5cm]D)--([yshift=0.5cm,xshift=0.5cm]D)--([yshift=-0.5cm,xshift=0.5cm]A)--([yshift=-0.5cm,xshift=-0.5cm]A)--([yshift=0.5cm,xshift=-0.5cm]A);


     \draw [aquablue,line width=2pt]([yshift=-0.3cm,xshift=0.3cm]A)-- ([yshift=-0.3cm,xshift=-0.3cm]B)--
     ([yshift=0.3cm,xshift=-0.3cm]B) --([yshift=0.3cm,xshift=0.3cm]B)--([yshift=-0.3cm,xshift=0.3cm]C)--([yshift=-0.3cm,xshift=-0.3cm]C)--([yshift=0.3cm,xshift=-0.3cm]C)-- ([yshift=0.3cm,xshift=0.3cm]D)--([yshift=-0.3cm,xshift=0.3cm]D)--([yshift=-0.3cm,xshift=-0.3cm]D)--([yshift=0.3cm,xshift=-0.3cm]A)--([yshift=0.3cm,xshift=0.3cm]A)--([yshift=-0.3cm,xshift=0.3cm]A);

\end{tikzpicture}
\begin{center}
    {\large \bfseries \color{aquablue!80!blue} \underline{Problem A : Size of Moon’s Shadow}}
\end{center}

Given,

\begin{itemize}
    \item Full eclipse began at $t_1=20:31:31$
    \item Full eclipse ended at $t_2= 20:32:32$
    \item  Moon’s shadow's speed $v=\qty{0.75}{km/s}$
\end{itemize}
$\therefore $ The time interval $\Delta t =t_2-t_1=00:01:01=\qty{61}{s}$\\
$\therefore$ The distance the Moon's shadow traveled,\\ $s=v\Delta t = \qty{0.75}{km/s}\cdot\qty{61}{s}=\qty{45.75}{km}$\\


\begin{tcolorbox}[
    enhanced,
    attach boxed title to top center={yshift=-3mm,yshifttext=-1mm},
    colback=black,
    colframe=aquamarine!50!aquablue,
    coltitle=blue!70!black,
    colbacktitle=aquablue,
    title=Answer,
    fonttitle=\bfseries,
    boxed title style={size=small,colframe=blue},
    coltext=ivory!85!aquablue
]
   \centerline{The width of the Moon’s shadow is \qty{45.75}{km}}
\end{tcolorbox}
\vspace{1cm}
\begin{center}
    {\large \bfseries \color{aquablue!80!blue} \underline{Problem B : Floating Box}}
\end{center}
Given,
\begin{itemize}
    \item Density of box $\rho_{box} = \qty{650}{kg/m^3}$ 
    \item Density of water $\rho_{water} = \qty{1000}{kg/m^3}$
    \item The box has outer length $a=\qty{25}{cm}$ , width $b=\qty{12}{cm}$ and hight $c=\qty{18}{cm}$
    \item The thickness of wall is $d=\qty{1.0}{cm}$
\end{itemize}
Two forces are acting on the box - wight and buoyancy. The net force is $\sum F =0$\\
$\therefore$ The volume of the box,\\ $V=abc-(a-2d)(b-2d)(c-2d)=\qty{5400}{cm^3}-\qty{3680}{cm^3}=\qty{1720}{cm^3}$\\
$\therefore$ The weight $W=mg=\rho_{box} V g$\\
We think that hight $h$ of the box is beneath the water.\\
$\therefore$ Volume of box under water $V_w=abh$\\
According to Archimedes law,\\
$\therefore$ The buoyant force $F_b=\rho_{water}V_wg=\rho_{water}abhg$\\
$\therefore\sum F= W-F_b=0 \implies F_b=W \implies \rho_{water}abhg=\rho_{box} V g$\\
$$\therefore h = \frac{\rho_{box} V }{\rho_{water}ab}=\frac{\qty{650}{kg/m^3}\cdot \qty{1720}{cm^3} }{\qty{1000}{kg/m^3}\cdot \qty{25}{cm} \cdot \qty{12}{cm}}\approx\qty{3.73}{cm}$$
$\therefore$ The top of the box above the water surface is $h_t=c-h\approx\qty{14.3}{cm}$

\begin{tcolorbox}[
    enhanced,
    attach boxed title to top center={yshift=-3mm,yshifttext=-1mm},
    colback=black,
    colframe=aquamarine!50!aquablue,
    coltitle=blue!70!black,
    colbacktitle=aquablue,
    title=Answer,
    fonttitle=\bfseries,
    boxed title style={size=small,colframe=blue},
    coltext=ivory!85!aquablue
]
  \centerline{The top of the box is about $\qty{14.3}{cm}$ high above the water surface.}
\end{tcolorbox}
\newpage
\begin{tikzpicture}[remember picture,overlay,inner sep=0,outer sep=0]
     \draw[blue,line width=4pt] ([xshift=-1.5cm,yshift=-2cm]current page.north east) coordinate (A)--([xshift=1.5cm,yshift=-2cm]current page.north west) coordinate(B)--([xshift=1.5cm,yshift=2cm]current page.south west) coordinate (C)--([xshift=-1.5cm,yshift=2cm]current page.south east) coordinate(D)--cycle;

     \draw [aquablue,line width=2pt]([yshift=0.5cm,xshift=-0.5cm]A)-- ([yshift=0.5cm,xshift=0.5cm]B)--
     ([yshift=-0.5cm,xshift=0.5cm]B) --([yshift=-0.5cm,xshift=-0.5cm]B)--([yshift=0.5cm,xshift=-0.5cm]C)--([yshift=0.5cm,xshift=0.5cm]C)--([yshift=-0.5cm,xshift=0.5cm]C)-- ([yshift=-0.5cm,xshift=-0.5cm]D)--([yshift=0.5cm,xshift=-0.5cm]D)--([yshift=0.5cm,xshift=0.5cm]D)--([yshift=-0.5cm,xshift=0.5cm]A)--([yshift=-0.5cm,xshift=-0.5cm]A)--([yshift=0.5cm,xshift=-0.5cm]A);


     \draw [aquablue,line width=2pt]([yshift=-0.3cm,xshift=0.3cm]A)-- ([yshift=-0.3cm,xshift=-0.3cm]B)--
     ([yshift=0.3cm,xshift=-0.3cm]B) --([yshift=0.3cm,xshift=0.3cm]B)--([yshift=-0.3cm,xshift=0.3cm]C)--([yshift=-0.3cm,xshift=-0.3cm]C)--([yshift=0.3cm,xshift=-0.3cm]C)-- ([yshift=0.3cm,xshift=0.3cm]D)--([yshift=-0.3cm,xshift=0.3cm]D)--([yshift=-0.3cm,xshift=-0.3cm]D)--([yshift=0.3cm,xshift=-0.3cm]A)--([yshift=0.3cm,xshift=0.3cm]A)--([yshift=-0.3cm,xshift=0.3cm]A);

\end{tikzpicture}

\begin{center}
    {\large \bfseries \color{aquablue!80!blue} \underline{Problem C : Echo Measure}}
\end{center}

Given,
\begin{itemize}
    \item The air temperature is $\qty{20}{\degree C}$
    \item The student hears an echo $t=\qty{1.24}{s}$
    \item For sound in air, use the approximate relation,
    $$v= \left( 331+0.60\unit{\frac{T}{\degree C}} \right) \unit{\frac{m}{s}}$$
\end{itemize}
\begin{center}
    {\large \bfseries \color{aquablue!80!blue} \underline{(a)}}
\end{center}
$\therefore$ The speed of sound,
$$v= \left( 331+0.60\unit{\frac{\qty{20}{\degree C}}{\degree C}} \right) \unit{\frac{m}{s}}=\qty{343}{m/s}$$\\
\begin{center}
    {\large \bfseries \color{aquablue!80!blue} \underline{(b)}}
\end{center}
We suppose that the cliff is distance $d$ away.\\
$\therefore$ The distance,
$$d=\frac{1}{2}vt=\frac{1}{2}\cdot\qty{343}{m/s}\cdot \qty{1.24}{s}\approx \qty{213}{m}$$\\
\begin{center}
    {\large \bfseries \color{aquablue!80!blue} \underline{(c)}}
\end{center}
Given, the new temperature $T=\qty{32}{\degree C}$\\
$\therefore$ The new speed of sound,
$$v_w= \left( 331+0.60\unit{\frac{\qty{32}{\degree C}}{\degree C}} \right) \unit{\frac{m}{s}}=\qty{350.2}{m/s}$$
$\therefore$ The echo retures in time,\\
$$t_w=\frac{2d}{v_w}=\frac{2\cdot \qty{213}{m}}{\qty{350.2}{m/s}}\approx \qty{1.22}{s}$$
$\therefore$ The echo returns earlier,
$$\Delta t =t_w-t=\qty{1.24}{s}-\qty{1.22}{s}=\qty{0.02}{s}$$\\

\begin{tcolorbox}[
    enhanced,
    attach boxed title to top center={yshift=-3mm,yshifttext=-1mm},
    colback=black,
    colframe=aquamarine!50!aquablue,
    coltitle=blue!70!black,
    colbacktitle=aquablue,
    title=Answer,
    fonttitle=\bfseries,
    boxed title style={size=small,colframe=blue},
    coltext=ivory!85!aquablue
]
  \centerline{(a) \qty{343}{m/s} ; (b) \qty{213}{m} ; (c) \qty{0.02}{s}}
\end{tcolorbox}
\newpage
\begin{tikzpicture}[remember picture,overlay,inner sep=0,outer sep=0]
     \draw[blue,line width=4pt] ([xshift=-1.5cm,yshift=-2cm]current page.north east) coordinate (A)--([xshift=1.5cm,yshift=-2cm]current page.north west) coordinate(B)--([xshift=1.5cm,yshift=2cm]current page.south west) coordinate (C)--([xshift=-1.5cm,yshift=2cm]current page.south east) coordinate(D)--cycle;

     \draw [aquablue,line width=2pt]([yshift=0.5cm,xshift=-0.5cm]A)-- ([yshift=0.5cm,xshift=0.5cm]B)--
     ([yshift=-0.5cm,xshift=0.5cm]B) --([yshift=-0.5cm,xshift=-0.5cm]B)--([yshift=0.5cm,xshift=-0.5cm]C)--([yshift=0.5cm,xshift=0.5cm]C)--([yshift=-0.5cm,xshift=0.5cm]C)-- ([yshift=-0.5cm,xshift=-0.5cm]D)--([yshift=0.5cm,xshift=-0.5cm]D)--([yshift=0.5cm,xshift=0.5cm]D)--([yshift=-0.5cm,xshift=0.5cm]A)--([yshift=-0.5cm,xshift=-0.5cm]A)--([yshift=0.5cm,xshift=-0.5cm]A);


     \draw [aquablue,line width=2pt]([yshift=-0.3cm,xshift=0.3cm]A)-- ([yshift=-0.3cm,xshift=-0.3cm]B)--
     ([yshift=0.3cm,xshift=-0.3cm]B) --([yshift=0.3cm,xshift=0.3cm]B)--([yshift=-0.3cm,xshift=0.3cm]C)--([yshift=-0.3cm,xshift=-0.3cm]C)--([yshift=0.3cm,xshift=-0.3cm]C)-- ([yshift=0.3cm,xshift=0.3cm]D)--([yshift=-0.3cm,xshift=0.3cm]D)--([yshift=-0.3cm,xshift=-0.3cm]D)--([yshift=0.3cm,xshift=-0.3cm]A)--([yshift=0.3cm,xshift=0.3cm]A)--([yshift=-0.3cm,xshift=0.3cm]A);

\end{tikzpicture}


\begin{center}
    {\large \bfseries \color{aquablue!80!blue} \underline{Problem D : Garden Pipe Split}}
\end{center}
\vspace{5pt}
Given,
 \begin{itemize}
    \item Splitter gauge pressure $\Delta p = 180\text{ kPa}$ 
    \item Both branches have length $L$ and cross-sectional area $A$.
    \item $\Delta p = kLv^2$
    \item $k_2(\lambda) = k_1(1 + 3\lambda^2)$\\
\end{itemize}

Substituting $v = Q/A$,
$$ \Delta p = \left(\frac{k L}{A^2}\right) Q^2 $$ \\

As $\Delta p_1=\Delta p_2$,
$$\left(\frac{k_1(1 + 3\lambda^2) L}{A^2}\right) Q_2^2=\left(\frac{k_1 L}{A^2}\right) Q_1^2$$
$$ Q_2 = \frac{Q_1}{\sqrt{1 + 3\lambda^2}}$$\\

By using conservation of mass at the splitter, the total flow rate $Q_{\text{tot}}$ is,
$$Q_{\text{tot}}(\lambda) = Q_1 + Q_2= Q_1 + \frac{Q_1}{\sqrt{1 + 3\lambda^2}}= Q_1 \left( 1 + \frac{1}{\sqrt{1 + 3\lambda^2}} \right)$$
$$Q_{\text{tot}}(\lambda) =Q_1 \left( 1 + \frac{1}{\sqrt{1 + 3\lambda^2}} \right) $$\\
Which is the required expression for the total flow rate through the splitter.
\begin{tcolorbox}[
    enhanced,
    attach boxed title to top center={yshift=-3mm,yshifttext=-1mm},
    colback=black,
    colframe=aquamarine!50!aquablue,
    coltitle=blue!70!black,
    colbacktitle=aquablue,
    title=Answer,
    fonttitle=\bfseries,
    boxed title style={size=small,colframe=blue},
    coltext=ivory!85!aquablue
]
  \centerline{It is shown that : $ Q_{\text{tot}}(\lambda) =Q_1 \left( 1 + \frac{1}{\sqrt{1 + 3\lambda^2}} \right)$}
\end{tcolorbox}
\newpage


\begin{center}
    {\large \bfseries \color{aquablue!80!blue} \underline{Problem E : Lightning Rod}}
\end{center}
During thunderstorm, charge separation creates a voltage \(V\) between the cloud and the ground, forming an electric field,
$$E = -\nabla V$$ 
The charge density \( \sigma \) accumulates densely on sharp points ,as like rod's tip, because :
$$\sigma \propto 1/R$$
Here,\\
\(R\) is the radius of curvature of the rod's tip . This high charge concentration increases the local electric field, as :
$$E = \sigma/\varepsilon_0$$
 The ionized air triggers a controlled electrical discharge path before the rest of the structure can. At the end, according to Ohm’s law,
 $$I = V/R$$
 The lightning current flows safely through the thick conductor directly into the Earth. Thus, our buildings remain safe due to these rods.

\begin{tcolorbox}[
    enhanced,
    attach boxed title to top center={yshift=-3mm,yshifttext=-1mm},
    colback=black,
    colframe=aquamarine!50!aquablue,
    coltitle=blue!70!black,
    colbacktitle=aquablue,
    title=Answer,
    fonttitle=\bfseries,
    boxed title style={size=small,colframe=blue},
    coltext=ivory!85!aquablue
]
  The rod concentrates charge at its tip to increase the local electric field which causes the lightning to strike the rod and ground through wire.
\end{tcolorbox}


\end{document}
